%==============================================================
\section{Các thuật toán lượng tử cơ bản}
\label{sec:algorithms}
%==============================================================

Gần như mọi thuật toán lượng tử cơ bản -- Deutsch--Jozsa, Bernstein--Vazirani, Grover, phase
estimation -- đều dùng chung một thủ thuật:
\begin{quote}
  \emph{Ghi câu trả lời của một hàm vào \textbf{phase} của amplitude, rồi dùng $H$ để biến phase thành
  bit đọc được.}
\end{quote}
Thủ thuật này gọi là \textbf{phase kickback}. Chương này xây dựng nó từ đầu, sau đó áp dụng lần lượt
vào các bài toán với độ tổng quát tăng dần, và kết thúc bằng hai giao thức dùng entanglement làm tài
nguyên chính.

\subsection{Oracle và phase kickback}

\paragraph{Oracle.} Cho hàm $f:\{0,1\}^n\to\{0,1\}$ mà ta chỉ được \emph{gọi} chứ không được xem bên
trong (giống gọi một API). Để đưa $f$ vào mạch lượng tử, nó phải được đóng gói thành một cổng unitary
-- và unitary thì phải khả nghịch. Phép $\ket{x}\mapsto\ket{f(x)}$ nói chung không khả nghịch, nên
cách chuẩn là thêm một qubit phụ (\textbf{ancilla}) và ghi kết quả bằng XOR:
\begin{equation}
  U_f\,\ket{x}\ket{y} = \ket{x}\ket{y\oplus f(x)}.
  \label{eq:oracle}
\end{equation}
Dạng này gọi là \textbf{bit-flip oracle} và luôn khả nghịch ($U_f^2 = I$).

\paragraph{Phase kickback.} Đặt ancilla ở $\ket{-}$ thay vì $\ket{0}$. Vì $\ket{-}$ là vector riêng
của $X$ với trị riêng $-1$, và oracle chính là ``áp $X$ lên ancilla khi $f(x)=1$'', ta có
\begin{equation}
  U_f\,\ket{x}\ket{-} = (-1)^{f(x)}\,\ket{x}\ket{-}.
  \label{eq:kickback}
\end{equation}
Ancilla không đổi; nó chỉ ``đá ngược'' (kick back) dấu $(-1)^{f(x)}$ vào qubit dữ liệu. Câu trả lời
của $f$ đã chuyển từ \emph{bit của ancilla} sang \emph{phase của qubit dữ liệu}. Khi qubit dữ liệu ở
superposition, một lần gọi oracle ghi dấu lên \emph{mọi} giá trị $x$ cùng lúc:
\begin{equation}
  U_f\,\ket{+}\ket{-} = \tfrac{1}{\sqrt2}\Big((-1)^{f(0)}\ket{0} + (-1)^{f(1)}\ket{1}\Big)\ket{-}.
\end{equation}

\subsection{Bài toán Deutsch và Deutsch--Jozsa}

\paragraph{Deutsch.} Cho $f:\{0,1\}\to\{0,1\}$, hỏi $f$ là \emph{hằng} ($f(0)=f(1)$) hay \emph{cân
bằng} ($f(0)\neq f(1)$)? Cổ điển phải gọi $f$ hai lần. Theo phương trình trên, qubit dữ liệu sau
oracle là $\pm\ket{+}$ nếu $f$ hằng và $\pm\ket{-}$ nếu $f$ cân bằng; áp $H$ rồi đo cho đúng
$f(0)\oplus f(1)$ -- chỉ với \textbf{một} lần gọi \cite{deutsch1985}.

\begin{figure}[h]
  \centering
  \begin{quantikz}
    \lstick{$x: \ket{0}$} & \gate{H} & & \gate[2]{U_f} & \gate{H} & \meter{} & \setwiretype{c} \rstick{$f(0)\oplus f(1)$} \\
    \lstick{$\text{anc}: \ket{0}$} & \gate{X} & \gate{H} & & & &
  \end{quantikz}
  \caption{Mạch Deutsch. Ancilla được đưa về $\ket{-}$ bằng $X$ rồi $H$; chỉ đo qubit dữ liệu.}
  \label{fig:deutsch}
\end{figure}

\paragraph{Deutsch--Jozsa.} Mở rộng cho $f:\{0,1\}^n\to\{0,1\}$ với lời hứa $f$ hoặc hằng, hoặc cân
bằng (đúng một nửa đầu vào cho 0) \cite{deutsch1992}. Cổ điển, trong trường hợp xấu nhất cần
$2^{n-1}+1$ lần gọi để chắc chắn. Mạch lượng tử giống hệt Deutsch với $n$ qubit dữ liệu:
$H^{\otimes n}$ -- oracle -- $H^{\otimes n}$ -- đo. Kết quả \textbf{toàn 0} nghĩa là $f$ hằng, bất
kỳ kết quả nào khác nghĩa là $f$ cân bằng. Lý do: amplitude của $\ket{0^n}$ sau cùng là
$\frac{1}{2^n}\sum_x(-1)^{f(x)}$, bằng $\pm1$ nếu $f$ hằng và bằng $0$ nếu $f$ cân bằng.

\begin{warnbox}
Chỉ đo qubit dữ liệu. Nếu dùng \code{measure\_all()}, ancilla (đang ở $\ket{-}$) cũng bị đo và cho bit
ngẫu nhiên, làm chuỗi counts có một bit ``rác'' khiến ta tưởng thuật toán sai.
\end{warnbox}

\subsection{Bernstein--Vazirani}

\paragraph{Bài toán.} Có chuỗi bí mật $s\in\{0,1\}^n$; oracle tính
$f(x) = s\cdot x \bmod 2 = s_0x_0\oplus\dots\oplus s_{n-1}x_{n-1}$. Tìm $s$ \cite{bernstein1997}.

Cổ điển, mỗi lần gọi (với $x$ là một vector đơn vị) chỉ lấy được \textbf{một} bit của $s$, nên cần
$n$ lần. Lượng tử dùng đúng mạch Deutsch--Jozsa. Oracle được hiện thực bằng một cổng CX từ qubit $i$
vào ancilla cho mỗi $i$ có $s_i=1$. Sau oracle, phase kickback cho trạng thái
\begin{equation}
  \frac{1}{\sqrt{2^n}}\sum_{x}(-1)^{s\cdot x}\ket{x} = H^{\otimes n}\ket{s}.
\end{equation}
Vì $H^{\otimes n}$ là nghịch đảo của chính nó, lớp $H^{\otimes n}$ cuối cùng đưa trạng thái về đúng
$\ket{s}$: đo ra $s$ với xác suất \textbf{100\%}, chỉ sau \textbf{một} lần gọi. Đây là ví dụ gọn nhất
về \emph{interference có chủ đích}: các amplitude sai triệt tiêu nhau, amplitude đúng cộng dồn lại.

\begin{keybox}[Lợi thế ở đây đo bằng gì?]
Thước đo là \textbf{số lần gọi oracle} (query complexity), không phải tổng số cổng: bên trong oracle
BV vẫn có tới $n$ cổng CX. Deutsch--Jozsa và BV là các bài toán được \emph{thiết kế} để lợi thế lộ ra;
chúng không nói lên rằng máy lượng tử nhanh hơn trong các bài toán thực tế.
\end{keybox}

\subsection{Thuật toán Grover}

Khác với hai bài toán trên, Grover giải một bài toán \textbf{hoàn toàn tổng quát}: tìm kiếm không cấu
trúc \cite{grover1996}. Có $N = 2^n$ khả năng, trong đó $M$ khả năng là ``nghiệm''; ta chỉ có oracle
trả lời ``đây có phải nghiệm không''. Cổ điển cần trung bình $N/2$ lần thử; Grover cần khoảng
$\frac{\pi}{4}\sqrt{N/M}$ lần -- lợi thế \textbf{bình phương}.

\paragraph{Phase oracle.} Đánh dấu nghiệm bằng dấu âm:
$U_\omega\ket{x} = -\ket{x}$ nếu $x=\omega$, ngược lại giữ nguyên. Hiện thực: áp $X$ lên các qubit
bằng 0 trong $\omega$, áp một cổng multi-controlled $Z$, rồi áp $X$ trả lại.

\paragraph{Diffuser.} Dấu âm một mình không đổi xác suất ($|-a|^2 = |a|^2$). Bước \textbf{diffuser}
$D = H^{\otimes n}(2\ket{0}\bra{0}-I)H^{\otimes n}$ thực hiện phép \emph{lật quanh giá trị trung
bình} (inversion about the mean): $a_k \mapsto 2\bar a - a_k$. Amplitude âm nằm dưới trung bình rất xa
nên sau khi lật nhảy lên trên trung bình -- tức được khuếch đại.

\begin{figure}[h]
  \centering
  \begin{quantikz}[row sep=0.25cm]
    \lstick[3]{$\ket{0}^{\otimes n}$} & \gate{H} & \gate[3]{U_\omega}\gategroup[3,steps=4,style={dashed,rounded corners,fill=AccentBlue!8,inner xsep=2pt},background,label style={label position=above,anchor=south,yshift=-0.1cm}]{lặp $k \approx \frac{\pi}{4}\sqrt{N/M}$ lần: oracle + diffuser} & \gate{H} & \gate[3]{2\ket{0}\bra{0}-I} & \gate{H} & \meter{} \\
    & \gate{H} & & \gate{H} & & \gate{H} & \meter{} \\
    & \gate{H} & & \gate{H} & & \gate{H} & \meter{}
  \end{quantikz}
  \caption{Cấu trúc thuật toán Grover: khởi tạo superposition đều, rồi lặp $k$ vòng gồm oracle
  $U_\omega$ và diffuser.}
  \label{fig:grover}
\end{figure}

\paragraph{Phân tích hình học.} Trạng thái luôn nằm trong mặt phẳng hai chiều sinh bởi ``tổng các
nghiệm'' và ``tổng các phi nghiệm''. Mỗi vòng Grover (oracle + diffuser) là một phép \emph{quay} góc
$2\theta$ về phía nghiệm, với $\sin\theta = \sqrt{M/N}$. Sau $k$ vòng,
\begin{equation}
  P(k) = \sin^2\big((2k+1)\theta\big), \qquad
  k_{\text{opt}} \approx \frac{\pi}{4}\sqrt{\frac{N}{M}}.
  \label{eq:grover}
\end{equation}
Công thức~\eqref{eq:grover} nói hai điều: số vòng tối ưu tăng như $O(\sqrt N)$, và \textbf{lặp nhiều
hơn không tốt hơn} -- quá $k_{\text{opt}}$, trạng thái quay vượt qua nghiệm và xác suất giảm
(Hình~\ref{fig:grover-prob}). Với $n=3$, $M=1$: $k_{\text{opt}}=2$ và $P\approx 94.5\%$, khớp với
kết quả mô phỏng trong notebook.

\begin{figure}[h]
  \centering
  \begin{tikzpicture}
    \begin{axis}[
      width=0.78\linewidth, height=5.2cm,
      xlabel={Số vòng Grover $k$}, ylabel={$P(k)$},
      xmin=-0.5, xmax=12.5, ymin=0, ymax=1.08,
      xtick={0,1,...,12}, ytick={0,0.25,0.5,0.75,1},
      grid=major, grid style={black!10},
      legend style={at={(0.5,1.03)}, anchor=south, legend columns=2, font=\footnotesize},
    ]
      \addplot[ybar, bar width=7pt, fill=AccentBlue!60, draw=AccentBlue,
               samples at={0,1,...,12}]
        {sin((2*x+1)*asin(0.25))^2};
      \addlegendentry{$N=16,\ M=1$}
      \addplot[only marks, mark=*, ProbRed, samples at={0,1,...,12}]
        {sin((2*x+1)*asin(sqrt(4/16)))^2};
      \addlegendentry{$N=16,\ M=4$}
      \draw[dashed, OkGreen, thick] (axis cs:3,0) -- (axis cs:3,1.05);
      \node[OkGreen, font=\footnotesize, anchor=west, fill=white, inner sep=1pt] at (axis cs:3.15,0.12) {$k_{\text{opt}}=3$};
    \end{axis}
  \end{tikzpicture}
  \caption{Xác suất tìm được nghiệm theo số vòng Grover (phương trình~\eqref{eq:grover}). Với một
  nghiệm trong 16 phần tử, xác suất đạt cực đại $96\%$ tại $k=3$ rồi giảm; với 4 nghiệm, số vòng
  tối ưu chỉ là 1.}
  \label{fig:grover-prob}
\end{figure}

\subsection{No-cloning, teleportation và superdense coding}

\paragraph{Định lý no-cloning.} Không tồn tại unitary nào sao chép được một trạng thái tuỳ ý:
$\ket{\psi}\ket{0}\mapsto\ket{\psi}\ket{\psi}$ với mọi $\psi$ \cite{wootters1982}. Chứng minh ngắn
dựa vào tính tuyến tính: nếu một mạch sao chép được $\ket{0}$ và $\ket{1}$ (CX làm được điều này),
thì tác dụng của nó lên $\ket{+}$ buộc phải là $\frac{1}{\sqrt2}(\ket{00}+\ket{11})$ -- một trạng thái
Bell -- chứ không phải $\ket{+}\ket{+}$.

\paragraph{Teleportation.} Alice muốn chuyển trạng thái chưa biết $\ket{\psi}$ của qubit $q_0$ cho Bob
\cite{bennett1993teleport}. Hai người chia sẵn một cặp Bell ($q_1$ của Alice, $q_2$ của Bob). Alice
áp CX$(q_0,q_1)$ rồi $H(q_0)$, đo hai qubit của mình được hai bit $m_0, m_1$, gửi chúng cho Bob qua
kênh cổ điển; Bob áp $X^{m_1}$ rồi $Z^{m_0}$ và thu được đúng $\ket{\psi}$. Bước sửa cuối dùng
\textbf{mid-circuit measurement} và điều khiển cổ điển (trong Qiskit 2.x: \code{with qc.if\_test(...)}).

\begin{figure}[h]
  \centering
  \begin{quantikz}[row sep=0.3cm]
    \lstick{$q_0: \ket{\psi}$} & & & \ctrl{1} & \gate{H} & \meter{} & \setwiretype{c} \rstick{$m_0$} \\
    \lstick{$q_1: \ket{0}$} & \gate{H} & \ctrl{1} & \targ{} & & \meter{} & \setwiretype{c} \rstick{$m_1$} \\
    \lstick{$q_2: \ket{0}$} & & \targ{} & & & & \gate{X^{m_1}} & \gate{Z^{m_0}} & \rstick{$\ket{\psi}$}
  \end{quantikz}
  \caption{Giao thức teleportation. Hai cổng đầu tạo cặp Bell giữa Alice ($q_1$) và Bob ($q_2$).
  Hai cổng sửa của Bob được điều khiển bởi hai bit cổ điển $m_0, m_1$.}
  \label{fig:teleport}
\end{figure}

Teleportation không vi phạm thuyết tương đối: trước khi nhận hai bit cổ điển, qubit của Bob ở trạng
thái hoàn toàn ngẫu nhiên (đo cho 50/50). Nó cũng không vi phạm no-cloning: trạng thái gốc của Alice
bị phá huỷ khi đo. Do mạch có phép đo giữa chừng, không thể kiểm tra bằng \code{Statevector}; cách
kiểm tra chuẩn là áp nghịch đảo của mạch chuẩn bị $\ket{\psi}$ lên qubit của Bob rồi đo -- phải được
\code{0} trong 100\% số shot.

\paragraph{Superdense coding.} Chiều ngược lại: gửi \textbf{2 bit cổ điển} bằng \textbf{1 qubit}
\cite{bennett1992dense}. Với một cặp Bell chia trước, Alice áp $I$, $X$, $Z$ hoặc $ZX$ lên qubit của
mình tuỳ theo hai bit cần gửi; bốn phép này biến cặp Bell thành bốn trạng thái Bell trực giao, nên Bob
phân biệt chắc chắn bằng cách đảo mạch tạo Bell (CX rồi $H$) và đo. Tổng số qubit được truyền vẫn là
2 (một qubit đi trước khi chia cặp), nên không vi phạm giới hạn Holevo.

\subsection{Tổng kết chương}

\begin{table}[h]
  \centering
  \caption{So sánh các thuật toán và giao thức đã trình bày.}
  \label{tab:algo-summary}
  \small
  \begin{tabularx}{\linewidth}{@{}lllX@{}}
    \toprule
    \textbf{Bài toán} & \textbf{Cổ điển} & \textbf{Lượng tử} & \textbf{Cơ chế chính} \\
    \midrule
    Deutsch & 2 lần gọi & 1 lần gọi & phase kickback + $H$ \\
    Deutsch--Jozsa & $2^{n-1}+1$ lần gọi & 1 lần gọi & phase kickback + $H^{\otimes n}$ \\
    Bernstein--Vazirani & $n$ lần gọi & 1 lần gọi & interference về đúng $\ket{s}$ \\
    Grover & $O(N)$ & $O(\sqrt{N})$ & oracle + diffuser (khuếch đại amplitude) \\
    Teleportation & -- & 1 qubit $=$ 1 cặp Bell + 2 bit & entanglement + điều khiển cổ điển \\
    Superdense coding & -- & 2 bit $=$ 1 cặp Bell + 1 qubit & 4 trạng thái Bell trực giao \\
    \bottomrule
  \end{tabularx}
\end{table}

Điểm chung về bản chất: các thuật toán này \textbf{không} tính $f$ nhanh hơn. Chúng trả lời một câu
hỏi \emph{toàn cục} về $f$ mà không cần biết từng giá trị $f(x)$, nhờ ghi thông tin vào phase và
dùng interference để đọc ra. Tất cả đều được viết dưới dạng mạch lý tưởng -- dùng cổng
multi-controlled, CX giữa hai qubit bất kỳ. Chương tiếp theo xem điều gì xảy ra khi đưa các mạch đó
lên phần cứng thật.
